% ============================================================================
% Chapter 20 --- Wave 2: Observability baseline
% ----------------------------------------------------------------------------
% Purpose : establish measurement before changing things --- a regression
%           detector for every subsequent wave. Collectors deployed in
%           shadow mode (D12.7), no commitment to a specific analytical
%           store at this stage.
% Page target: ~5 pages.
% Dependencies: Wave 0 (Ch. 18); reference layer Ch. 12.
% Mirrors wave template (D18.2).
% ----------------------------------------------------------------------------
% Decisions registered in _coherence-EN.md:
%   D20.1 Wave 2 is purely additive --- existing logging untouched
%   D20.2 collectors deployed in shadow mode before any product commitment
%   D20.3 the canonical TAC schema is enforced at the collector layer
%         even when no TAC is yet emitted
%   D20.4 exit criteria reference Observability conformance subset (D16.5)
% ============================================================================

\chapter{Wave~2 --- Observability baseline}
\label{ch:wave2-observability}

\begin{caveatbox}[Dependencies]
\noindent Wave~2 applies the wave chapter template established in
\trchap{ch:wave0-foundations}{}. Its pre-requisite is Wave~0
(governance and inventory in place); it can proceed in parallel
with Wave~1. It instantiates the observability contract of
\trchap{ch:observability}{}, including the canonical TAC schema of
\S\ref{sec:ob-schema}, and its closure is gated by the four
Observability tests of \trchap{ch:conformance}{}.
\end{caveatbox}

Wave~2 is the wave that measures the institution before any
substantial change is made to it. The reasoning is procedural: an
institution that begins technical transformation without a baseline
of its current state has no objective way to recognise regression.
It is also the lowest-risk wave of the playbook --- the wave is
purely additive, and the existing logging arrangements continue
unchanged for as long as the institution chooses.

%-----------------------------------------------------------------------------
\section{Goal and dependencies}
\label{sec:w2-goal}

The goal of Wave~2 is to deploy a collector layer that captures
events from every layer of the architecture in the canonical TAC
schema of \trchap{ch:observability}{} (\S\ref{sec:ob-schema}), to
route those events to a portable ingestion pipeline, and to
establish the retention schedule that the data classification
warrants. The wave commits to the collector layer and the
ingestion pipeline; it does not commit to a specific analytical
store, which is treated as a substitutable consumer.

\paragraph{Why before other waves.} Subsequent waves change
operational state. Without a baseline, the change cannot be
distinguished from drift, regression, or pre-existing condition.
Wave~2's purpose is to make those distinctions possible.

%-----------------------------------------------------------------------------
\section{Pre-requisites}
\label{sec:w2-prereq}

Beyond Wave~0, five operational conditions strengthen Wave~2:

\begin{itemize}[leftmargin=1.5em,nosep]
  \item the log sources of the existing infrastructure are
    catalogued, with their volume, format, and current retention
    documented;
  \item the institution's data classification scheme (or the
    intent to develop one) is articulated, since the retention
    schedule is keyed on it;
  \item the existing SIEM (or the explicit absence thereof) is
    acknowledged, with its contractual envelope documented;
  \item the security operations function and its staffing profile
    --- whether delivered by an internal team, by a contracted
    managed-security provider, or by a distributed arrangement
    shared across units --- is identified, since the wave will
    produce events that the function will consume;
  \item where the adopter's labour law gives staff-representation
    bodies consultation or co-determination rights over workplace
    monitoring, the consultation of \loc{staff-representation-body}
    on the observability design (scope, metadata-only default,
    retention, access) is initiated --- deployment must not precede
    it where the law so requires.
\end{itemize}

%-----------------------------------------------------------------------------
\section{Coexistence pattern}
\label{sec:w2-coexistence}

Wave~2 introduces no displacement. The existing log paths continue
to operate; existing dashboards, alerts, and detection rules
remain in place; the existing analytical store, if any, continues
to receive its incumbent feed.

\paragraph{Collector layer in shadow mode.} The new collectors are
deployed at the event-producing sources and forward to the new
ingestion pipeline (an open-source log-routing agent feeding an
open-source distributed event-streaming bus). Initially, the
events flow only into a development or
staging analytical store, where the institution exercises the
canonical schema, the retention policy, and the correlation
patterns. The production analytical store --- whether it remains
the existing one or is migrated later --- consumes from the
ingestion pipeline only when the institution has decided which
analytical store it wishes to commit to (a decision that is
neither blocked nor pre-empted by Wave~2).

\paragraph{Operating two layers.} For the duration of the wave, the
institution operates the legacy log paths and the new collector
layer simultaneously. The cost is real but bounded; the benefit is
that any subsequent wave's regression is observable against the
baseline established here.

%-----------------------------------------------------------------------------
\section{Trajectory: greenfield}
\label{sec:w2-greenfield}

For institutions in a greenfield position, Wave~2 establishes the
collector layer, the ingestion pipeline, and the analytical store
together. There is no incumbent to coexist with, and the wave is
correspondingly compressed (indicative duration four to six
months). The greenfield trajectory has the freedom to commit to an
analytical store at wave open, since the substitution risk is
absent.

%-----------------------------------------------------------------------------
\section{Trajectory: brownfield single-suite-centric}
\label{sec:w2-brownfield-ms}

For institutions whose existing observability is dominated by
a cloud-vendor-native SIEM (for example Microsoft Sentinel with
Log Analytics), Wave~2 introduces an open-source collector layer,
built on a vendor-neutral telemetry-collection standard, upstream
of the incumbent SIEM, without displacing its role.

\paragraph{Pattern.} Open-source collectors, such as those built on
a vendor-neutral telemetry-collection standard, are deployed at
every event-producing source (identity, network, endpoint,
compute). The events are converted to the canonical TAC
schema and forwarded to two destinations: the incumbent SIEM (which
continues its role for the operational SOC) and a
parallel open-source store (an open-source search-and-analytics
engine or an open-source log-aggregation store) for portability and
for the retention horizons that the incumbent SIEM's pricing makes
unattractive. The institution retains the incumbent SIEM for the
analytical-depth use cases and acquires the data-portability
property of the architecture's contract.

\paragraph{What changes in Wave~2.} The institution gains a
canonical schema across its sources, an ingestion pipeline whose
re-routing is tactical rather than structural, and a retention
horizon for cold-tier events that does not pay the per-GB premium
of the production SIEM. The collector deployment is itself
conformance-tested before the wave closes.

\paragraph{What does not change.} The incumbent SIEM remains the
production SOC tool, with its dashboards, detections, and the team's
familiarity. The incumbent vendor's licensing trajectory is
unaffected by this wave; the institution gains the option to revisit it in
future waves without operational disruption.

\paragraph{Indicative duration.} Six to twelve months, dominated by
the source-by-source deployment of collectors and by the
canonical-schema enforcement effort.

%-----------------------------------------------------------------------------
\section{Trajectory: brownfield hybrid}
\label{sec:w2-brownfield-hybrid}

For institutions whose existing observability is built on Splunk
or a comparable per-volume-priced analytical store, Wave~2
introduces the same collector layer upstream of Splunk, with the
same parallel open-source store, for analogous reasons.

\paragraph{Pattern.} Open-source collectors built on a
vendor-neutral telemetry-collection standard deployed at every
event-producing source, with conversion to the canonical TAC
schema. The events flow to Splunk (incumbent SOC store) and to a
parallel open-source search-and-analytics engine or log-aggregation
store. Sampling decisions become
tactical at the collector layer (D12.8: sample by TAC session, not
by event), reducing the per-volume exposure to Splunk without
losing the activity that the institution has chosen to make
visible.

\paragraph{What changes in Wave~2.} The volume reaching Splunk is
brought under deliberate control; the parallel open store
preserves the events that volume management would otherwise have
discarded; the canonical schema is enforced uniformly, removing
the schema drift that ad hoc Splunk-side parsing accumulates over
years.

\paragraph{What does not change.} Splunk remains the production
analytical surface. The detection content (alerts, correlations,
dashboards) is preserved.

\paragraph{Indicative duration.} Six to twelve months.

%-----------------------------------------------------------------------------
\section{Reversibility criteria}
\label{sec:w2-reversibility}

Wave~2 is structurally reversible: removing the collectors
returns the institution to its prior state without operational
impact. Three explicit conditions support this:

\begin{itemize}[leftmargin=1.5em,nosep]
  \item the legacy log paths are not modified, so a removal of the
    new collectors leaves them operational;
  \item the analytical-store decision is deferred to a subsequent
    wave or wave-internal milestone, so no commitment is made
    that would be costly to reverse;
  \item the parallel open-source store, if exercised during the
    wave, is sized for the wave's purpose and can be decommissioned
    without contractual penalty.
\end{itemize}

%-----------------------------------------------------------------------------
\section{Exit criteria}
\label{sec:w2-exit}

Wave~2 closes when the four Observability tests of the conformance
suite of \trchap{ch:conformance}{} (\S\ref{sec:cf-suite}) pass on
the new collector layer, when the agreed source coverage is
achieved, and when the retention schedule is documented and
verified.

\begin{itemize}[leftmargin=1.5em,nosep]
  \item Events from every layer arrive in the canonical TAC schema
    at the parallel store; the retention horizon is verified by
    deliberate retrieval at the documented limit; cross-layer
    correlation via \texttt{correlation\_id} succeeds on a
    representative test activation; event re-routing to an
    alternative analytical store succeeds on a sampling window.
    (Observability conformance subset.)
  \item Source coverage: the institutionally agreed proportion of
    event-producing sources is feeding the new collector layer.
    The remaining sources are either in a documented schedule or
    in the documented exception register.
  \item Retention schedule: published, with the documented horizons
    per data class, and verified at least once during the wave.
\end{itemize}

\begin{realworldbox}[Audit pipeline precedents]
\noindent Several research-intensive institutions are reported to
operate observability pipelines combining open-source collectors
with commercial analytical stores. Specific architectures
discussed in the fora of \loc{nren-name} and \loc{nren-csirt}, and
in the security operations environments of large research
infrastructures, would benefit from being consulted in their
current published form before being cited; the patterns described
in this chapter are consistent with those reported in those
settings.
\end{realworldbox}

%-----------------------------------------------------------------------------
\section{Regulatory anchoring}
\label{sec:w2-regulatory}

\noindent Wave~2 establishes the observability and audit
baseline. It is a key contributor to \gls{gdpr}~Art.~30 (records
of processing, derivable from the canonical TAC schema), Art.~32
(security of processing, monitored through measurable
indicators) and Art.~33--34 (breach notification, supported by
the timeline-reconstruction substrate produced by the wave). It
also supports the incident-handling obligations of the
applicable national cybersecurity baseline regime --- the
national transposition of the NIS2 Directive where it applies ---
under its incident-handling provision (NIS2~Art.~21(2)(b)) and
prepares the ground for that regime's incident-reporting
timelines (NIS2~Art.~23), fully exercised at Wave~7. The wave
is the principal anchor of
27001~A.8.15 (logging) and A.8.16 (monitoring activities) and
of NIST CSF \textsf{DETECT} (DE.CM, DE.AE). Detailed mapping in
\trchap{ch:regulatory-mapping}{} \S\ref{sec:rm-gdpr},
\S\ref{sec:rm-nis2}, \S\ref{sec:rm-iso},
\S\ref{sec:rm-nist-csf}.
